% Section 1.7, translated from sv/chapters/01/exercises.tex.

\section{Exercises}
\label{c1:sec:uppgifter}

\begin{aside}
\textbf{How to work with the exercises.} Parts~1 and~2 are done during the lesson: first on your
own, a few minutes per exercise, then together. Part~3 is extra exercises for further practice
after the lesson. Most of them can be done in your head or with pencil and paper.

\facitnotis{L01}
\end{aside}

% ------------------------------------------------------------------------------------------
\uppgiftsdel{Part 1}{Introductory exercises}

\uppgift{1.1}{Calculations with the order of operations}
Calculate the following expressions \textbf{without} a calculator.
\begin{delar}(4)
  \task $3 + 2 \times 5$
  \task $(3 + 2) \times 5$
  \task $12 - 4 \times 2 + 1$
  \task $\dfrac{10 + 2}{3} \times 4 - 1$
\end{delar}

\uppgift{1.2}{Fractions}
Calculate and simplify.
\begin{delar}(4)
  \task $\dfrac{1}{4} + \dfrac{1}{6}$
  \task $\dfrac{3}{5} - \dfrac{1}{4}$
  \task $\dfrac{2}{3} \times \dfrac{9}{4}$
  \task $\dfrac{5}{6} \div \dfrac{5}{12}$
\end{delar}

\uppgift{1.3}{Percentages}
\begin{parts}
  \item Calculate $40\,\%$ of $250$.
  \item What percentage of $72$ is $18$?
  \item A resistor has the nominal value $100\,\Omega$ with a tolerance of $\pm 5\,\%$. Give the
    lowest and the highest permitted resistance.
\end{parts}

% ------------------------------------------------------------------------------------------
\uppgiftsdel{Part 2}{New material}

\uppgift{2.1}{Resistors in parallel}
Two resistors $R_1 = 12\,\Omega$ and $R_2 = 6\,\Omega$ are connected in parallel. The total
resistance $R_{\text{TOT}}$ is given by
\[
  \frac{1}{R_{\text{TOT}}} = \frac{1}{R_1} + \frac{1}{R_2}.
\]
\begin{delar}(2)
  \task Calculate $\dfrac{1}{R_{\text{TOT}}}$.
  \task Calculate $R_{\text{TOT}}$.
\end{delar}

\uppgift{2.2}{Series and parallel connection}
The circuit below consists of the resistor $R_1$ in series with the parallel connection of $R_2$
and $R_3$.

\bild[0.5\linewidth]{lectures/L01/appendix/images/circuit_series_parallel.png}

The component values are $R_1 = 4\,\text{k}\Omega$, $R_2 = 24\,\text{k}\Omega$ and
$R_3 = 8\,\text{k}\Omega$. The parallel resistance $R_{\text{p}} = R_2 \parallell R_3$ can be
calculated in two ways:
\[
  \frac{1}{R_{\text{p}}} = \frac{1}{R_2} + \frac{1}{R_3}
  \qquad \text{or} \qquad
  R_{\text{p}} = \frac{R_2 \times R_3}{R_2 + R_3}
\]
\begin{parts}
  \item Calculate $R_{\text{p}}$ with both formulas and check that they give the same answer.
  \item Calculate the circuit's total resistance $R_{\text{TOT}} = R_1 + R_{\text{p}}$.
  \item Why must the parallel connection be calculated before $R_1$ is added? Relate your answer
    to the order of operations.
\end{parts}

\uppgift{2.3}{Current and voltage in the circuit}
The circuit in \uppref{2.2} is supplied with the voltage $U = 20\,\text{V}$. Ohm's law relates
voltage, resistance and current:
\[
  U = R \times I
\]

\begin{aside}
\note[Tip] If you calculate with the resistance in k$\Omega$ and the voltage in V, you get the
current directly in mA.
\end{aside}

\begin{parts}
  \item Calculate the current $I$ that the voltage source delivers.
  \item Calculate the voltage $U_1$ across $R_1$.
  \item Calculate the voltage $U_{\text{p}}$ across the parallel connection.
  \item Calculate the currents $I_2$ and $I_3$ through $R_2$ and $R_3$ respectively. Check that
    $I_2 + I_3 = I$.
\end{parts}

\uppgift{2.4}{Absolute value}
An AC voltage varies between $-8\,\text{V}$ and $+8\,\text{V}$.
\begin{parts}
  \item What is the amplitude $\abs{U}$ of the voltage?
  \item At a certain instant, the voltage is $u(t) = -5.3\,\text{V}$. Calculate $\abs{u(t)}$.
\end{parts}

\uppgift{2.5}{Efficiency}
A battery-powered motor takes in $P_{\text{in}} = 24\,\text{W}$ and delivers
$P_{\text{out}} = 18\,\text{W}$.
\begin{parts}
  \item Calculate the motor's efficiency $\eta$ as a percentage.
  \item If the efficiency is improved to $90\,\%$ at the same input power, what will the output
    power be?
\end{parts}

\uppgift{2.6}{Number sets}
Place each number in the \emph{smallest} number set it belongs to: $\mathbb{N}$, $\mathbb{Z}$,
$\mathbb{Q}$, the irrational numbers or $\mathbb{R}$.
\begin{delar}(6)
  \task $7$
  \task $-3$
  \task $\dfrac{3}{4}$
  \task $\sqrt{2}$
  \task $0.25$
  \task $\pi$
\end{delar}

% ------------------------------------------------------------------------------------------
\uppgiftsdel{Part 3}{Extra exercises}
The exercises below are extra practice, best done on your own after the lesson.

\uppgift{3.1}{Order of operations with negative numbers}
Calculate \textbf{without} a calculator.
\begin{delar}(3)
  \task $-3 + 5 \times (-2)$
  \task $(-4)^2 - 4^2$
  \task $\dfrac{-12}{-3} + 2 \times (-5)$
  \task $8 - 2(3 - 7)$
  \task $\dfrac{(-2)^3 + 10}{2}$
\end{delar}

\uppgift{3.2}{The number line and absolute value}
\begin{parts}
  \item Sort the numbers in order of size, smallest first: $-2.5$, $\abs{-3}$, $\dfrac{7}{2}$,
    $-\abs{4}$, $0$.
  \item Calculate $\abs{-7} - \abs{3}$.
  \item Calculate $\abs{3 - 8}$ and $\abs{8 - 3}$. What does the result say about the distance
    between two numbers?
  \item Two measured values are $U_1 = 4.8\,\text{V}$ and $U_2 = 5.1\,\text{V}$. Write the size of
    the difference with an absolute value and calculate it.
\end{parts}

\uppgift{3.3}{Fractions, decimals and percentages}
Fill in the empty cells. Give the fractions in their simplest form.

\begin{narrowtable}{@{}*{3}{>{\centering\arraybackslash}p{6em}}@{}}
  \thead{Fraction} & \thead{Decimal} & \thead{Percentage} \\
  \midrule
  $\dfrac{1}{4}$ & & \\
  & $0.6$ & \\
  & & $12.5\,\%$ \\
  $\dfrac{3}{8}$ & & \\
\end{narrowtable}

\uppgift{3.4}{Fractions in several steps}
Calculate and simplify.
\begin{delar}(4)
  \task $\dfrac{2}{3} + \dfrac{1}{6} - \dfrac{1}{2}$
  \task $\dfrac{3}{4} \times \dfrac{8}{9} \div \dfrac{2}{3}$
  \task $\dfrac{\frac{1}{2} + \frac{1}{3}}{\frac{5}{6}}$
  \task $\dfrac{1}{\frac{1}{4} + \frac{1}{4}}$
  \task* Which circuit calculation do you recognise in d)?
\end{delar}

\uppgift{3.5}{Percentage change}
\begin{parts}
  \item A voltage rises from $12\,\text{V}$ to $15\,\text{V}$. By how many per cent does it rise?
  \item A current falls from $400\,\text{mA}$ to $340\,\text{mA}$. By how many per cent does it
    fall?
  \item A resistance first increases by $10\,\%$ and then decreases by $10\,\%$. Is it back at its
    original value? Justify your answer.
  \item A power of $60\,\text{W}$ is increased by $25\,\%$. What is the new power?
\end{parts}

\uppgift{3.6}{Tolerance and measurement deviation}
\begin{parts}
  \item A resistor is marked $4.7\,\text{k}\Omega \pm 10\,\%$. Give the lowest and the highest
    permitted value.
  \item The resistor measures $4.9\,\text{k}\Omega$. Is it within tolerance?
  \item By how many per cent does the measured value deviate from the nominal value?
  \item Another resistor is marked $220\,\Omega \pm 5\,\%$ and measures $208\,\Omega$. Does it
    pass?
\end{parts}

\uppgift{3.7}{Number sets: true or false}
Decide whether the statement is true or false, and justify your answer briefly.
\begin{parts}
  \item Every integer is also a rational number.
  \item Every rational number is also an integer.
  \item $\sqrt{9}$ is an irrational number.
  \item The sum of two irrational numbers is always irrational.
  \item The decimal $0.333\ldots$ is a rational number.
  \item Every real number can be written as a fraction of two integers.
\end{parts}

\uppgift{3.8}{Equal resistors in parallel}
\begin{parts}
  \item Two resistors of $100\,\Omega$ each are connected in parallel. Calculate $R_{\text{p}}$.
  \item Four resistors of $100\,\Omega$ each are connected in parallel. Calculate $R_{\text{p}}$.
  \item Show in general that $n$ equal resistors $R$ in parallel give
    $R_{\text{p}} = \dfrac{R}{n}$.
  \item How many $100\,\Omega$ resistors must be connected in parallel to give $20\,\Omega$?
\end{parts}

\uppgift{3.9}{Find the unknown resistor}
\begin{parts}
  \item $R_1 = 30\,\Omega$ is connected in parallel with $R_2$ and gives $R_{\text{p}} =
    12\,\Omega$. Calculate $R_2$.
  \item $R_1 = 220\,\Omega$ is to be supplemented with a resistor in series so that
    $R_{\text{TOT}} = 500\,\Omega$. Calculate that resistor.
  \item Check in a) that $R_{\text{p}}$ is smaller than the smallest of the resistors in it.
  \item Can a parallel connection of $R_1 = 30\,\Omega$ and some other resistor give
    $R_{\text{p}} = 40\,\Omega$? Justify your answer.
\end{parts}

\uppgift{3.10}{Ohm's law: fill in the table}
Fill in the empty cells using Ohm's law $U = R \times I$.

\begin{aside}
\note[Tip] With the resistance in k$\Omega$ and the voltage in V, the current comes out directly in
mA.
\end{aside}

\begin{narrowtable}{@{}p{5em}p{5em}p{5em}@{}}
  \thead{$U$} & \thead{$R$} & \thead{$I$} \\
  \midrule
  $12\,\text{V}$ & $4\,\text{k}\Omega$ & \\
  & $2.2\,\text{k}\Omega$ & $5\,\text{mA}$ \\
  $9\,\text{V}$ & & $3\,\text{mA}$ \\
  $5\,\text{V}$ & $1\,\text{k}\Omega$ & \\
\end{narrowtable}

\uppgift{3.11}{Efficiency in two stages}
A power supply consists of a transformer with the efficiency $\eta_1 = 95\,\%$ followed by a
rectifier with the efficiency $\eta_2 = 80\,\%$. The input power is $P_{\text{in}} = 200\,\text{W}$.
\begin{parts}
  \item Calculate the power after the transformer.
  \item Calculate the output power after the rectifier.
  \item Calculate the total efficiency $\eta_{\text{TOT}}$.
  \item How much power is lost as heat in total?
\end{parts}
